\tableofcontents


% =========================================================
% ГЛАВА 1. СБЛИЖЕНИЕ С СОХРАНЕНИЕМ СУММЫ
% =========================================================

\chapter{Сближение с сохранением суммы}

Метод Штурма начинается с очень простой операции: мы берём две переменные и
делаем их ближе друг к другу, сохраняя некоторую величину неизменной.
Сначала разберём самый важный случай: сохраняется сумма.

\begin{taskbox}{Сближение с сохранением суммы}{blue}{}
Пусть $a,b\ge 0$ и
\[
    a+b=S,
\]
где $S$ --- константа. Пусть, не умаляя общности, $a<b$.

Сближение с сохранением суммы --- это замена
\[
    (a,b)\longmapsto (a+\varepsilon,\; b-\varepsilon).
\]
Сумма действительно не меняется:
\[
    (a+\varepsilon)+(b-\varepsilon)=a+b=S.
\]

Чтобы точки не перескочили через середину отрезка $[a,b]$, нужно брать
\[
    0<\varepsilon\le \frac{b-a}{2}.
\]
Если
\[
    \varepsilon=\frac{b-a}{2},
\]
то обе переменные становятся равными:
\[
    a+\varepsilon=b-\varepsilon=\frac{a+b}{2}.
\]
Если же хотим, чтобы они остались строго по разные стороны от середины, то
берём
\[
    0<\varepsilon<\frac{b-a}{2}.
\]
\end{taskbox}

\begin{center}
    \begin{asy}
    import graph;
    size(13cm,3.6cm);
    
    defaultpen(fontsize(11pt));
    pen axispen = black+1.0bp;
    pen movePen = blue+1.1bp;
    pen dashedPen = gray+0.8bp+dashed;
    
    // координаты точек на оси
    real A = 0;
    real Ap = 1.6;
    real M = 4.2;
    real Bp = 6.8;
    real B = 8.4;
    
    // ось
    draw((A-0.8,0)--(B+0.6,0), axispen, Arrow(TeXHead));
    
    // засечки
    void tick(real x, string lab, pair shift=0)
    {
        draw((x,-0.12)--(x,0.12), axispen);
        label(lab, (x,0)+shift);
    }
    
    tick(A, "$a$", (0,-0.42));
    tick(Ap, "$a+\varepsilon$", (0,0.42));
    tick(M, "$\dfrac{a+b}{2}$", (0.4,-0.3));
    tick(Bp, "$b-\varepsilon$", (0,0.42));
    tick(B, "$b$", (0,-0.42));
    
    // середина
    draw((M,-0.95)--(M,1.0), dashedPen);
    label("середина", (M,1.15));
    
    // стрелки сближения
    draw(A..(0.45,0.55)..(1.0,0.55)..Ap, movePen, Arrow(TeXHead));
    label("$+\varepsilon$", (0.8,0.78), movePen);
    
    draw(B..(7.95,0.55)..(7.35,0.55)..Bp, movePen, Arrow(TeXHead));
    label("$-\varepsilon$", (7.55,0.78), movePen);
    
    // подпись
    draw((Ap,0)--(Bp,0), movePen);
    label("\textcolor{blue}{новые точки ближе друг к другу}", ((Ap+Bp)/2,-0.82));
    
    \end{asy}
\end{center}

\begin{taskbox}{Идея метода Штурма}{blue}{}
    Главная мысль в том, что мы смотрим, как при такой замене меняются разные выражения. Если нужная сторона неравенства меняется только в одну сторону, то можно постепенно сближать переменные, пока они не станут равными.
\end{taskbox}

\section{Расстояние}

Расстояние между числами это ни что иное как модуль их разность $|a-b|$. Понятное дело, что если мы <<сближаем числа>>, то расстояние между ними уменьшается. Строго это доказывается так:

\[
b'-a'=(b-\varepsilon)-(a+\varepsilon)=b-a-2\varepsilon.
\]
Так как $\varepsilon>0$, то
\[
b'-a'<b-a.
\]

\medskip

\section{Произведение}
После замены $(a,b)\longmapsto (a+\varepsilon,\; b-\varepsilon).$ получим

\[
a'b'=(a+\varepsilon)(b-\varepsilon)
=ab+\varepsilon(b-a)-\varepsilon^2.
\]
Поэтому
\[
a'b'-ab=\varepsilon(b-a)-\varepsilon^2
=\varepsilon(b-a-\varepsilon)>0.
\]

Последнее неравенство верно, так как мы условились, что $\varepsilon \leqslant \dfrac{b-a}{2}$. Значит, $a'b'>ab,$ то есть при сближении произведение увеличивается. Это отлично сходится с известной логикой, что \textbf{при фиксированном периметре прямоугольник имеет наибольшую площадь тогда и только тогда когда он квадрат.}
\medskip

\section{Сумма квадратов}

В случае с суммой квадратов нам даже не обязательно делать дурацкую замену с дурацким эпсилон. У нас уже достаточно знаний! Действительно,

$$
a^2 + b^2 = (a+b)^2 - 2ab
$$

Квадрат суммы остается незименным при сближении, а вот произведение как мы уже выяснили увеличивается. Но так как перед ним знак минус, значит сумма квадратов будет уменьшаться.

\medskip

\section{Сумма обратных величин}

Аналогично предыдущему случаю используем уже полученные знания,

$$
\dfrac{1}{a} + \dfrac{1}{b} = \dfrac{a+b}{ab}
$$

При сближении числитель не меняется, а вот знаменатель увеличивается. Значит сама дробь уменьшается.

\medskip

\section{Произведение $(1+a)(1+b)$}

Имеем
\[
(1+a)(1+b)=1+a+b+ab.
\]
При сближении сумма $a+b$ не меняется, а произведение $ab$ увеличивается. Поэтому сама величина увеличивается

\begin{remarkbox}
    При отдалении переменных все эффекты становятся противоположными: расстояние увеличивается, произведение уменьшается, сумма квадратов увеличивается, сумма обратных величин увеличивается, а $(1+a)(1+b)$ уменьшается.    
\end{remarkbox}


\medskip

\section{Таблица}



\begin{center}
\small
\setlength{\tabcolsep}{4pt}
\renewcommand{\arraystretch}{1.35}

\begin{center}
\small
\setlength{\tabcolsep}{6pt}
\renewcommand{\arraystretch}{1.35}

\begin{tabularx}{\textwidth}{
    >{\centering\arraybackslash}p{0.28\textwidth}
    >{\centering\arraybackslash}X
    >{\centering\arraybackslash}X
}
\toprule
\textbf{Величина}
&
\textbf{При сближении переменных}
&
\textbf{При отдалении переменных}
\\
\midrule

$a+b$
&
не меняется
&
не меняется
\\

\addlinespace[2mm]

$|a-b|$
&
уменьшается
&
увеличивается
\\

\addlinespace[2mm]

$ab$
&
увеличивается
&
уменьшается
\\

\addlinespace[2mm]

$a^2+b^2$
&
уменьшается
&
увеличивается
\\

\addlinespace[2mm]

$\dfrac1a+\dfrac1b,\quad a,b>0$
&
уменьшается
&
увеличивается
\\

\addlinespace[2mm]

$(1+a)(1+b)$
&
увеличивается
&
уменьшается
\\

\bottomrule
\end{tabularx}
\end{center}
\normalsize

\end{center}
\normalsize

\newpage

\chapter{Сближение с сохранением произведения}

Теперь разберём второй важный тип сближения: сохраняется произведение.

\begin{taskbox}{Сближение с сохранением произведения}{blue}{}
Пусть $a,b>0$ и
\[
    ab=P,
\]
где $P$ --- константа. Пусть, не умаляя общности, $a<b$.

Сближение с сохранением произведения --- это замена
\[
    (a,b)\longmapsto \left(a\varepsilon,\; \frac{b}{\varepsilon}\right).
\]
Произведение действительно не меняется:
\[
    (a\varepsilon)\cdot \frac{b}{\varepsilon}=ab=P.
\]

Чтобы первая переменная увеличивалась, а вторая уменьшалась, берём
\[
    \varepsilon>1.
\]
Чтобы переменные не перескочили через точку равенства, нужно брать
\[
    a\varepsilon\le \frac{b}{\varepsilon}.
\]
Отсюда
\[
    \varepsilon^2\le \frac ba,
\]
то есть
\[
    1<\varepsilon\le \sqrt{\frac ba}.
\]

Если
\[
    \varepsilon=\sqrt{\frac ba},
\]
то обе переменные становятся равными:
\[
    a\varepsilon=\frac b\varepsilon=\sqrt{ab}.
\]
Если же хотим, чтобы они остались строго по разные стороны от точки равенства, то берём
\[
    1<\varepsilon<\sqrt{\frac ba}.
\]
\end{taskbox}

\begin{center}
    \begin{asy}
    import graph;
    size(13cm,3.6cm);
    
    defaultpen(fontsize(11pt));
    pen axispen = black+1.0bp;
    pen movePen = blue+1.1bp;
    pen dashedPen = gray+0.8bp+dashed;
    
    // координаты точек на оси
    real A = 0;
    real Ap = 1.8;
    real M = 4.2;
    real Bp = 6.6;
    real B = 8.4;
    
    // ось
    draw((A-0.8,0)--(B+0.6,0), axispen, Arrow(TeXHead));
    
    // засечки
    void tick(real x, string lab, pair shift=0)
    {
        draw((x,-0.12)--(x,0.12), axispen);
        label(lab, (x,0)+shift);
    }
    
    tick(A, "$a$", (0,-0.42));
    tick(Ap, "$a\varepsilon$", (0,0.42));
    tick(M, "$\sqrt{ab}$", (0,-0.42));
    tick(Bp, "$\dfrac{b}{\varepsilon}$", (0,0.42));
    tick(B, "$b$", (0,-0.42));
    
    // точка равенства
    draw((M,-0.95)--(M,1.0), dashedPen);
    label("точка равенства", (M,1.15));
    
    // стрелки сближения
    draw(A..(0.45,0.55)..(1.1,0.55)..Ap, movePen, Arrow(TeXHead));
    label("$\cdot \varepsilon$", (0.9,0.78), movePen);
    
    draw(B..(7.95,0.55)..(7.25,0.55)..Bp, movePen, Arrow(TeXHead));
    label("$:\varepsilon$", (7.45,0.78), movePen);
    
    // подпись
    draw((Ap,0)--(Bp,0), movePen);
    label("\textcolor{blue}{новые точки ближе друг к другу}", ((Ap+Bp)/2,-0.82));
    
    \end{asy}
\end{center}

\begin{taskbox}{Идея}{blue}{}
    Теперь мы смотрим, как меняются разные выражения при замене
    \[
        (a,b)\longmapsto \left(a\varepsilon,\; \frac{b}{\varepsilon}\right),
    \]
    где произведение $ab$ остаётся неизменным.
\end{taskbox}

\section{Отношение большего к меньшему}

Так как $a<b$, отношение большего к меньшему равно
\[
    \frac ba.
\]
После замены получаем
\[
    \frac{b'}{a'}
    =
    \frac{\frac b\varepsilon}{a\varepsilon}
    =
    \frac{b}{a\varepsilon^2}.
\]
Так как $\varepsilon>1$, то
\[
    \frac{b}{a\varepsilon^2}<\frac ba.
\]
Значит, при сближении отношение большего к меньшему уменьшается.

\medskip

\section{Сумма}

После замены получаем
\[
    a'+b'=a\varepsilon+\frac b\varepsilon.
\]
Сравним это со старой суммой:
\[
    a'+b'-(a+b)
    =
    a\varepsilon+\frac b\varepsilon-a-b.
\]
Преобразуем:
\[
    a'+b'-(a+b)
    =
    a(\varepsilon-1)-b\left(1-\frac1\varepsilon\right) = (\varepsilon-1)\left(a-\frac b\varepsilon\right).
\]

Но из условия
\[
    1 \leq \varepsilon\le \sqrt{\frac ba}
\]
следует, что первая скобка положительная, а вторая отрицательная. Значит,

Значит,
\[
    a'+b'-(a+b)<0.
\]
Следовательно,
\[
    a'+b'<a+b.
\]
То есть при сближении с сохранением произведения сумма уменьшается.

\medskip

\section{Расстояние}

После замены расстояние между переменными равно
\[
    b'-a'=\frac b\varepsilon-a\varepsilon.
\]
Сравним старое расстояние с новым:
\[
    (b-a)-(b'-a')
    =
    b-a-\frac b\varepsilon+a\varepsilon.
\]
Преобразуем:
\[
    (b-a)-(b'-a')
    =
    b\left(1-\frac1\varepsilon\right)+a(\varepsilon-1).
\]
Так как $\varepsilon>1$, то обе части положительны:
\[
    b\left(1-\frac1\varepsilon\right)>0,
    \qquad
    a(\varepsilon-1)>0.
\]
Значит,
\[
    (b-a)-(b'-a')>0.
\]
Следовательно,
\[
    b'-a'<b-a.
\]
То есть при сближении расстояние уменьшается.

\medskip

\section{Сумма квадратов}

Используем уже полученное знание о сумме:
\[
    a^2+b^2=(a+b)^2-2ab.
\]
При сближении с сохранением произведения величина $ab$ не меняется, а сумма $a+b$ уменьшается. Поэтому квадрат суммы тоже уменьшается:
\[
    (a'+b')^2<(a+b)^2.
\]
Значит,
\[
    (a')^2+(b')^2<a^2+b^2.
\]
То есть сумма квадратов уменьшается.

\medskip

\section{Сумма обратных величин}

Имеем
\[
    \frac1a+\frac1b=\frac{a+b}{ab}.
\]
При сближении с сохранением произведения знаменатель $ab$ не меняется, а числитель $a+b$ уменьшается. Поэтому сама дробь уменьшается:
\[
    \frac1{a'}+\frac1{b'}<\frac1a+\frac1b.
\]

\medskip

\section{Произведение $(1+a)(1+b)$}

Имеем
\[
    (1+a)(1+b)=1+a+b+ab.
\]
При сближении с сохранением произведения величина $ab$ не меняется, а сумма $a+b$ уменьшается. Поэтому
\[
    (1+a')(1+b')<(1+a)(1+b).
\]
Значит, величина $(1+a)(1+b)$ уменьшается.

\newpage

\begin{remarkbox}
    При отдалении переменных все эффекты становятся противоположными: отношение большего к меньшему увеличивается, сумма увеличивается, расстояние увеличивается, сумма квадратов увеличивается, сумма обратных величин увеличивается, а $(1+a)(1+b)$ увеличивается. Произведение $ab$ при этом не меняется.
\end{remarkbox}

\medskip

\section{Таблица}

\begin{center}
\small
\setlength{\tabcolsep}{6pt}
\renewcommand{\arraystretch}{1.35}

\begin{tabularx}{\textwidth}{
    >{\centering\arraybackslash}p{0.28\textwidth}
    >{\centering\arraybackslash}X
    >{\centering\arraybackslash}X
}
\toprule
\textbf{Величина}
&
\textbf{При сближении переменных}
&
\textbf{При отдалении переменных}
\\
\midrule

$ab$
&
не меняется
&
не меняется
\\

\addlinespace[2mm]

$\dfrac ba,\quad 0<a<b$
&
уменьшается
&
увеличивается
\\

\addlinespace[2mm]

$a+b$
&
уменьшается
&
увеличивается
\\

\addlinespace[2mm]

$|a-b|$
&
уменьшается
&
увеличивается
\\

\addlinespace[2mm]

$a^2+b^2$
&
уменьшается
&
увеличивается
\\

\addlinespace[2mm]

$\dfrac1a+\dfrac1b,\quad a,b>0$
&
уменьшается
&
увеличивается
\\

\addlinespace[2mm]

$(1+a)(1+b)$
&
уменьшается
&
увеличивается
\\

\bottomrule
\end{tabularx}
\end{center}
\normalsize

\chapter{Алгоритм метода Штурма}

\begin{enumerate}[label=\textbf{Шаг \arabic*.}]
    \item Найти, какая величина в условии должна сохраняться: сумма, произведение, сумма квадратов, сумма углов.
    \item Выбрать две переменные $x,y$ и заменить их на более близкие $x',y'$ с тем же ограничением.
    \item Доказать, что нужная часть неравенства меняется монотонно: увеличивается или уменьшается.
    \item Повторять сближение. Обычно пределом будет $x_1=\ldots=x_n$; иногда край, где одна переменная равна нулю; для трёх переменных часто остаётся случай $x=y$.
    \item Проверить предельный случай. Он обычно превращается в одно переменное неравенство.
\end{enumerate}

\begin{tcolorbox}[
    enhanced,
    breakable,
    colback=gray!10,
    colframe=blue!70!black,
    colbacktitle=blue!85!black,
    coltitle=white,
    fonttitle=\bfseries,
    title=Мини-словарь,
    boxrule=0.8pt,
    arc=2mm
]
\textbf{Сблизить} две переменные --- заменить их на более равные.\par
\textbf{Раздвинуть} две переменные --- заменить их на более далёкие, часто доводя одну до края области.\par
\textbf{Край} --- ситуация, где одна переменная равна $0$, или центральный угол стремится к $0$, или другая естественная граница.
\end{tcolorbox}

\chapter{Как методом Штурма уничтожается неравенство о средних}

Разберём цепочку средних для положительных чисел $x_1,\ldots,x_n$:
\[
\frac{n}{\frac1{x_1}+\cdots+\frac1{x_n}}
\le
\sqrt[n]{x_1\cdots x_n}
\le
\frac{x_1+\cdots+x_n}{n}
\le
\sqrt{\frac{x_1^2+\cdots+x_n^2}{n}}.
\]
Слово ``уничтожается'' здесь означает: неравенство сводится к случаю, когда все переменные равны.

\section{Геометрическое среднее не больше арифметического}

Пусть сумма $S=x_1+\cdots+x_n$ фиксирована. Сблизим пару $x_i,x_j$:
\[
    x_i+x_j=s,\qquad x_i x_j\le\left(\frac{s}{2}\right)^2.
\]
При замене пары на $s/2,s/2$ произведение всех переменных не уменьшается, а сумма остаётся той же. Повторяя операцию, приходим к равным переменным $S/n$. Значит, при фиксированной сумме произведение максимально при равенстве:
\[
    x_1\cdots x_n\le\left(\frac{S}{n}\right)^n.
\]
Берём корень:
\[
    \sqrt[n]{x_1\cdots x_n}\le\frac{x_1+\cdots+x_n}{n}.
\]

\section{Гармоническое среднее не больше геометрического}

Пусть произведение $P=x_1\cdots x_n$ фиксировано. Сблизим пару $x_i,x_j$, сохраняя произведение:
\[
    x_i x_j=p,
    \qquad
    x_i,x_j\longrightarrow \sqrt p,\sqrt p.
\]
Тогда
\[
    \frac1{x_i}+\frac1{x_j}
    =\frac{x_i+x_j}{p}
    \ge\frac{2\sqrt p}{p}
    =\frac{2}{\sqrt p}.
\]
То есть при сближении сумма обратных чисел не увеличивается. Повторяем, пока все числа не станут равны $\sqrt[n]{P}$. Тогда
\[
    \frac1{x_1}+\cdots+\frac1{x_n}\ge \frac{n}{\sqrt[n]{P}},
\]
а значит
\[
    \frac{n}{\frac1{x_1}+\cdots+\frac1{x_n}}
    \le \sqrt[n]{P}.
\]

\section{Арифметическое среднее не больше квадратического}

Пусть сумма $S$ фиксирована. При сближении пары $x,y$ с суммой $s$ величина
\[
    x^2+y^2=(x+y)^2-2xy
\]
уменьшается, потому что $xy$ увеличивается. Значит, при фиксированной сумме сумма квадратов минимальна при равенстве:
\[
    x_1^2+\cdots+x_n^2\ge n\left(\frac{S}{n}\right)^2.
\]
Отсюда
\[
    \frac{x_1+\cdots+x_n}{n}\le
    \sqrt{\frac{x_1^2+\cdots+x_n^2}{n}}.
\]

\begin{remarkbox}
Обратите внимание: мы почти ничего не считали. Метод сам сказал, где искать худший случай: там, где переменные равны. После этого остаётся подставить равные значения.
\end{remarkbox}

\chapter{Разбор некоторых задач}

\phantomsection
\addcontentsline{toc}{section}{Задача 1}
\begin{taskbox}{Задача 1}{blue}{}
Докажите, что
\[
    (1+x_1)(1+x_2)\cdots(1+x_n)\ge 2^n,
\]
где $x_1x_2\cdots x_n=1$ и $x_i>0$.
\end{taskbox}

\textbf{Решение методом Штурма.}
Сохраняем произведение. Возьмём пару $x,y$ и заменим её на $\sqrt{xy},\sqrt{xy}$. Тогда произведение всех переменных не изменится. А вот сама величина в задаче уменьшится (см. таблицу 2.7)

Повторяя сближение, мы можем только уменьшать левую часть, пока все переменные не станут равными. Из условия $x_1\cdots x_n=1$ в равном случае получаем
\[
    x_1=\cdots=x_n=1.
\]
Значит, минимальное значение левой части равно
\[
    (1+1)^n=2^n.
\]
Неравенство доказано. Равенство достигается только при $x_1=\cdots=x_n=1$.

\phantomsection
\addcontentsline{toc}{section}{Задача 2}
\begin{taskbox}{Задача 2}{blue}{}
Пусть сумма положительных чисел $x_1,x_2,\ldots,x_n$ равна $1$. Докажите, что
\[
    \frac{(1-x_1)(1-x_2)\cdots(1-x_n)}{x_1x_2\cdots x_n}
    \ge (n-1)^n.
\]
\end{taskbox}

\textbf{Решение методом Штурма.}
Сохраняем сумму. Возьмём пару $x,y$ и обозначим $x+y=s$. Остальные переменные не трогаем. Парный множитель равен
\[
    \frac{(1-x)(1-y)}{xy}
    =\frac{1-s+xy}{xy}
    =1+\frac{1-s}{xy}.
\]
Так как вся сумма равна $1$, имеем $s\le1$, причём $1-s\ge0$. При сближении пары с фиксированной суммой произведение $xy$ увеличивается, поэтому величина
\[
    1+\frac{1-s}{xy}
\]
не увеличивается. Значит, сближение не увеличивает всё выражение.

Повторяя операцию, уменьшаем левую часть до случая
\[
    x_1=\cdots=x_n=\frac1n.
\]
Тогда
\[
    \frac{(1-\frac1n)^n}{(\frac1n)^n}=(n-1)^n.
\]
Значит, исходное выражение не меньше $(n-1)^n$.

\phantomsection
\addcontentsline{toc}{section}{Задача 3. IMO 1984}
\begin{taskbox}{Задача 3. IMO 1984}{blue}{}
Пусть $x,y,z\ge0$ и $x+y+z=1$. Докажите, что
\[
    0\le xy+yz+zx-2xyz\le \frac7{27}.
\]
Если $x,y,z>0$, то левая часть строго положительна.
\end{taskbox}

\textbf{Левая оценка.}
Так как $0\le x,y,z\le1$, то
\[
    xy\ge xyz,
    \qquad yz\ge xyz,
    \qquad zx\ge xyz.
\]
Следовательно,
\[
    xy+yz+zx-2xyz\ge xyz\ge0.
\]

\textbf{Правая оценка методом Штурма.}
Обозначим
\[
    F(x,y,z)=xy+yz+zx-2xyz.
\]
Пусть $x\le y\le z$. Тогда $y\le\frac12$. Сблизим две крайние переменные $x$ и $z$, сохранив их сумму; переменную $y$ оставим фиксированной. При фиксированных $x+z$ и $y$ имеем
\[
    F= y(x+z)+xz-2xyz=y(x+z)+xz(1-2y).
\]
Так как $1-2y\ge0$, то при сближении пары $x,z$ произведение $xz$ увеличивается, а значит $F$ не уменьшается. Следовательно, максимум достигается в случае, когда две переменные равны.

Положим $x=z=a$, $y=1-2a$, где $0\le a\le\frac12$. Тогда
\[
    F(a,1-2a,a)=2a(1-2a)+a^2-2a^2(1-2a)
    =4a^3-5a^2+2a.
\]
Остаётся проверить
\[
    4a^3-5a^2+2a\le \frac7{27},\qquad 0\le a\le\frac12.
\]
Производная равна
\[
    12a^2-10a+2=2(2a-1)(3a-1).
\]
Значит, максимум на отрезке достигается при $a=\frac13$ или на концах. Подстановка даёт
\[
    F\left(\frac13,\frac13,\frac13\right)=\frac7{27},
    \qquad
    F(0,1,0)=0,
    \qquad
    F\left(\frac12,0,\frac12\right)=\frac14.
\]
Так как $\frac14<\frac7{27}$, максимум равен $\frac7{27}$.


\phantomsection
\addcontentsline{toc}{section}{Задача 4.}
\begin{taskbox}{Задача 4}{blue}{}
Пусть $a,b,c\ge0$ и $a+b+c=1$. Докажите, что
\[
    a^3+b^3+c^3+6abc\ge \frac14.
\]
\end{taskbox}

\textbf{Решение.}
Пусть $a\ge b\ge c$. Тогда $c\le\frac13$. Зафиксируем $c$ и сумму $a+b=1-c$. Получаем
\[
\begin{aligned}
    a^3+b^3+c^3+6abc
    &=(a+b)^3-3ab(a+b)+c^3+6abc\\
    &=(1-c)^3+c^3-3ab(1-3c).
\end{aligned}
\]
Так как $1-3c\ge0$, выражение минимально тогда, когда $ab$ максимально. При фиксированной сумме $a+b$ произведение $ab$ максимально при $a=b$. Значит, достаточно рассмотреть
\[
    a=b=\frac{1-c}{2},\qquad 0\le c\le\frac13.
\]
Тогда
\[
\begin{aligned}
    a^3+b^3+c^3+6abc
    &=2\left(\frac{1-c}{2}\right)^3+c^3+6\left(\frac{1-c}{2}\right)^2c\\
    &=\frac{9c^3-9c^2+3c+1}{4}.
\end{aligned}
\]
Остаётся заметить, что
\[
    9c^3-9c^2+3c=3c(3c^2-3c+1)\ge0,
\]
поскольку $3c^2-3c+1>0$ для всех $c$. Следовательно, выражение не меньше $\frac14$.

Равенство достигается при $(a,b,c)=(\frac12,\frac12,0)$ и перестановках.


\phantomsection
\addcontentsline{toc}{section}{Задача 5. IMO 1995}
\begin{taskbox}{Задача 5. IMO 1995}{blue}{}
Пусть $a,b,c>0$ и $abc=1$. Докажите, что
\[
    \frac{1}{a^3(b+c)}+
    \frac{1}{b^3(c+a)}+
    \frac{1}{c^3(a+b)}\ge\frac32.
\]
\end{taskbox}

\textbf{Решение.}
Когда в условии стоит произведение $abc=1$, полезно искать замену, которая выравнивает переменные. Положим
\[
    x=\frac1a,
    \qquad y=\frac1b,
    \qquad z=\frac1c.
\]
Тогда $xyz=1$, и
\[
    \frac{1}{a^3(b+c)}=\frac{x^3}{\frac1y+\frac1z}
    =\frac{x^3yz}{y+z}
    =\frac{x^2}{y+z},
\]
потому что $xyz=1$. Поэтому надо доказать
\[
    \frac{x^2}{y+z}+\frac{y^2}{z+x}+\frac{z^2}{x+y}\ge\frac32.
\]
По неравенству Коши
\[
    \frac{x^2}{y+z}+\frac{y^2}{z+x}+\frac{z^2}{x+y}
    \ge
    \frac{(x+y+z)^2}{2(x+y+z)}
    =\frac{x+y+z}{2}.
\]
А из $xyz=1$ по AM-GM имеем
\[
    x+y+z\ge3.
\]
Следовательно, вся сумма не меньше $\frac32$.

\begin{remarkbox}
Это решение не является чистым сближением, но мысль та же: условие $abc=1$ говорит, что естественная точка равенства --- $a=b=c=1$. Замена $x=1/a$ делает выражение подходящим для стандартного неравенства.
\end{remarkbox}